\documentclass[bachelor, och, assignment]{SCWorks}
\usepackage[T2A]{fontenc}
\usepackage[utf8]{inputenc}
\usepackage{graphicx}
\usepackage[sort,compress]{cite}
\usepackage{amsmath}
\usepackage{amssymb}
\usepackage{amsthm}
\usepackage{fancyvrb}
\usepackage{longtable}
\usepackage{array}
\usepackage[english,russian]{babel}
\usepackage[colorlinks=true]{hyperref}

\begin{document}

% Кафедра (в родительном падеже)
\chair{дискретной математики и информационных технологий}

% Тема работы
\title{Разработка системы стеганографического сокрытия данных в сетевых пакетах и модуля обнаружения вредоносных вложений}

% Курс
\course{4}

% Группа
\group{421}

% Специальность/направление код - наименование
\napravlenie{09.03.01 "--- Информатика и вычислительная техника}

% Фамилия, имя, отчество в родительном падеже
\author{Иванова Ивана Ивановича}

% Заведующий кафедрой
\chtitle{доцент, к.\,ф.-м.\,н.}
\chname{Л.\,Б.\,Тяпаев}

% Научный руководитель
\satitle{канд. техн. наук, доцент}
\saname{И.\,О.\,Фамилия}

% Семестр
\term{8}

% Наименование практики
\practtype{Преддипломная}

% Продолжительность практики
\duration{4}

% Даты практики
\practStart{03.02.2026}
\practFinish{11.06.2026}

% Год выполнения
\date{2026}

\maketitle

Исследование проблемы обнаружения скрытых каналов передачи информации в сетевых протоколах TCP/IP является актуальной задачей информационной безопасности. Анализ деятельности APT-группировок (Turla, Equation Group, APT29) показывает активное использование стеганографических методов сокрытия данных в заголовках и полезной нагрузке сетевых пакетов для управления вредоносным ПО и эксфильтрации конфиденциальных данных. Существующие системы защиты (DPI, IDS/IPS, DLP) не обеспечивают надёжного обнаружения таких каналов.

Целью работы является разработка программной системы стеганографического сокрытия данных в сетевых пакетах с криптографической защитой и многоуровневым модулем обнаружения скрытых каналов и вредоносных вложений.

Для достижения поставленной цели необходимо решить следующие задачи:

1. Провести анализ существующих методов сетевой стеганографии и классификацию скрытых каналов в протоколах TCP/IP.

2. Реализовать стеганографические каналы в пяти полях протоколов: IP~ID, TCP~ISN, TCP~Timestamp, ICMP~Payload и DNS~QNAME с криптографической защитой AES-256-GCM.

3. Разработать протокол фрагментации NST с двухуровневым контролем целостности (CRC32, SHA-256) и поддержкой неупорядоченной доставки чанков.

4. Разработать модуль обнаружения скрытых каналов, объединяющий статистический анализ (хи-квадрат, KS-тест, энтропия Шеннона), алгоритмы машинного обучения (Isolation Forest, Random Forest) и сигнатурный поиск.

5. Реализовать модуль анализа содержимого payload для обнаружения вредоносных вложений (20~сигнатур в 7~категориях).

6. Провести экспериментальное тестирование системы и оценить эффективность стеганографических каналов и модуля обнаружения.

\textbf{Основные результаты:}

Разработана система NetStego, реализующая полный конвейер стеганографической передачи данных через пять каналов TCP/IP с пропускной способностью от 20~байт/с до 115{,}9~Кбит/с. Модуль обнаружения обеспечивает F1-меру от 0{,}80 до 0{,}92 при уровне ложных срабатываний не более 10\,\%. Система включает 269~автоматизированных тестов (134~модульных, 68~интеграционных, 67~тестов целостности).
\vspace{\baselineskip}


{\bf Срок предоставления работы:} 11.06.2026
\vspace{\baselineskip}

Рассмотрено на заседании кафедры дискретной математики и

информационных технологий
\vspace{\baselineskip}

Протокол № 3 от 10.10.2025
\vspace{\baselineskip}

Секретарь  \quad  \underline{\hspace{4cm}}   \quad Кузнецов Т.М.
\vspace{\baselineskip}

Дата выдачи задания 10.10.2025
\vspace{\baselineskip}

Задание получил \quad  \underline{\hspace{4cm}}   \quad Иванов И.И.


\end{document}
